
\section{Terms of Trade and the Effects of Trade}

\begin{frame}{Winners and Losers from Trade}
    \showmytikz[1,2,3,4,5]{../shared/tikz/tikz-freight}
    \mycomment{Let us re-visit the entire trade value chain; there are winners, and potentially losers, in both countries, and along the route.}
\end{frame}


\begin{frame}{Who Wins Form Trade?}
    \mycomment{
        The advantages of trade for overall welfare are easily proven. However, large advantages of consumers contrast with sometimes very sharp disadvantages of (certain) producers, as trade may out-price local producers altogether.
    
        Often, producers are better organised and have a stronger lobby than consumers. Also, if producers are locally concentrated, shifts in production may impact local communities, as factories and entire industries close.
    
        Thus unsurprisingly, producers often lobby the state to intervene, while consumers rarely possess the same lobbying strength.
    }
    \mytext{
        \begin{itemize}
            \item Consumers 
            \item Producers of the comparatively cheaper good
            \item Trades with a link to those
            \item Traders, logistic companies, ports etc.
            \item The state via customs, tariffs and taxes (?)
        \end{itemize}
    }
\end{frame}

\begin{frame}{Who Wins Form Trade?}

    
    \showmybitmap{../shared/bitmaps/wto-value-added-global-trade.png}
    \cite{wto2023}

    \mycomment{
    According to WTO statistics, importers gain more of total value added than exporters. This would indicate that consumer surpluses outsize producer surpluses, which is unsurprising.
    }
    
    
\end{frame}
   


\begin{frame}{Gains for the Traders}
\mycomment{
The theory of comparative advantage does not look at the merchants and traders, who also play a not insignificant role. Often the trade industry is overlooked.
}
\mytext{
    Traders can win:
    \begin{itemize}
        \item as long as prices are different: arbitrage
        \item as long as trade is a monopoly: monopoly rent
        \item after that: marginal pricing (no profit)
        \item Trader should just be earning from logistics
        \item speculative wins
    \end{itemize}
}
\end{frame}

\begin{frame}{Gains for the Traders}
\mycomment{Let us re-visit the different parties along any trading route, from upstream to downstream:}
\mytext{
    \begin{itemize}
        \item producers
        \item exporters and importers
        \item of commodities, of itermediate, of manufactured goods
        \item of services and intellectual properties
        \item of capital; investors
        \item insurers, financial services
        \item local logistic providers
        \item retailers
        \item consumers
        \item states?
    \end{itemize}
}
\end{frame}


\begin{frame}{Income and Substitution Effect}
    \mycomment{
    We remember from Microeconomics that any relative price change has income and substitution effects.
    }
    \showmytikz[1,2,3,4]{../shared/tikz/tikz-consumer}
\end{frame}

\begin{frame}{Long Term Effects}
    \mytext{
        \begin{itemize}
            \item prices of tradable goods go down
            \item (relative) prices of non-tradable goods go up
            \item Technological change?
            \item Work force size / migration?
            \item Movement of capital?
            \item Convergence, divergence?
            \item Economies of scale?
            \item Rising production and consumption globally
        \end{itemize}
    }
\end{frame}



\begin{frame}{Terms of Trade}
    \mytext{
    Countries import and export different (bundles of) goods, for a number of reasons.
    
    ToT is the relation of the price of the exported good(s) to the imported good(s):

        \begin{equation}
            ToT=\frac{P_{Ex}}{P_{Im}}
        \end{equation}

        When speaking at the ToT of a country, we typically analyse a bundle of goods (the main import and export goods in the trade balance).

        This links to the exchange rate issue: imports and exports are typically denominated in different currencies.
    }

        \mycomment{
            For now, we assume the Terms of Trade be real exchange rates. There is, however, an established link between 
            changes in the exchange rate and the ToT. A depreciation of the home currency can affect imported and exported goods differently, but the overall effect on the Current Account depend on the elasticities of both bundles of goods           
            (\noun{Marshall-Learner}-Condition).
        }
    
\end{frame}

\begin{frame}{Terms of Trade}
    \mytext{
        The price of primary commodities declines relative to that of manufactured goods over the long term (\noun{Prebisch-Singer} hypothesis)

        The production of manufactured goods requires skill and capital, and shows larger economies of scale. 

        Income elasticity of demand and of labour required is different for primary and manufactured goods.
        
        Specialising on primary commodities may be a development hindering dead end road. \cite{galeano1973}
    }
\end{frame}

\begin{frame}{Terms of Trade}
    \mytext{
        The elasticities of both demand and of supply of commodities are inelastic in the short term.

        As a result, prices fluctuate more than those of manufactured goods.
    }
\end{frame}


\begin{frame}{Terms of Trade}
\mytext{
        Prebisch-Singer thesis:\\
        ToT worsen for exporters of commodities\\
        due to different income elasticity of demand


        \begin{itemize}
            \item Import substitution industrialisation (ISI) 1950ies-1980ies
            \item Export-oriented industrialisation (ESI) today; model: East~Asia (`tigers')
        
        \end{itemize}
    }
\end{frame}

\begin{frame}{Non-Tradable Goods}
    \mytext{
    Some goods are non-tradable:
    \begin{itemize}
        \item personal services
        \item labour?
    \end{itemize}
    That explains why a hair cut is vastly cheaper in some countries than in others, while tradable goods cost almost the same across borders.
    }
    \mycomment{
        Economies have a different endowment with labour. The labour force is mainly staying put: work migration is negligible in comparison to the trade of goods.

        This results in stark differences between the costs of labour.
    }
\end{frame}


\begin{frame}{Dutch Desease}
\mytext{
    \begin{itemize}
        \item booming sector (new gas field
        \item appreciates currency
        \item `lagging' sector disadvantage
        \item inflation
        \item problem to stabilise exchange rate
    \end{itemize}
}
\end{frame}

\begin{frame}{Rentier State}
\mycomment{
    Follow the money, goes the old adage, if you want to understand motives of people and reasons for action. Some states gain most of their revenue from a single industry - typically, mining or agriculture - rather than from income taxes. There is then not much incentive to invest in productivity growth. As long as that work-free income can be redistributed to keep everyone happy, this may be a good model indeed. 

    However, many such states seem lacking in democracy and participation.     
}
    \mytext{
       \begin{itemize}
        \item most government revenue from one sector
        \item independent of productivity
        \item redistribution of wealth
        \item unaccountable elites
        \item petrostates (Iran, Gabon...)
        \item democracy?
    \end{itemize} 
    }
\end{frame}

\begin{frame}{Foreign Direct Investment FDI}
\mycomment{
    Rather than importing a good into a country, a company may decide to invest directly into that country and produce right there. 
    
    This would (to some degree) substitute transactions in the trade balance (as part of the current account) by transactions in the capital and financial account. 

    Foreign direct investment comes with much longer commitment than trade, which can be shifted pretty much instantaneously, and is thus much more linked to macroeconomic considerations.
}
    \mytext{
     \begin{itemize}
        \item buy or build 
        \item non-financial (non-portfolio)
        \item country risk
        \item substitution, complementary
        \item tax competition
        \item diversification, risk minimisation?
        \item market access?
        \item repatriation?
        \item mergers \& acquisitions?
    \end{itemize}   
    }
\end{frame}



\begin{frame}{What About the Losers of Trade?}
\mycomment{
    Theory goes a long way to explain the overall welfare gains of trade for consumers and producers. However, the losers of trade deserve at least some consideration.
}
\mytext{
    \begin{itemize}
        \item Producers of the comparatively more expensive good (in a country)
        \item Trades with a link to those
        \item Environment?
        \item Trade balance?
        \item Local traditions and trades?
    \end{itemize}
}
\end{frame}

\begin{frame}{What About the Losers of Trade?}
\mycomment{
    A good question is whether the losers or the winners of trade are better organised, and more able to push their agenda to politics. Consumers, the main benefactors of trade, are typically less organised than producers, where production is often geographically concentrated. Free trade wipes out entire industries, leading to all sorts of protectionism.
}
\mytext{
    \begin{itemize}
        \item sectorial and local organisations
        \item lobbying 
        \item subsidies, protectionism
        \item consumers are less organised
    \end{itemize}
    Kaldor-Hicks criterium: can the gains be used to compensate for the losses?
}
\end{frame}


\begin{frame}{Gini Coefficient and Lorenz Curve}
    \showmytikz{../shared/tikz/tikz-gini}
    \mycomment{

    A related question is how trade influences not only overall income, but inequality in a country.

    Gini coefficient of 0 means total equality\\
    Gini coefficient of 1 means total inequality


    \begin{itemize}
        \item problematic in many aspects
        \item how about demographic?
        \item how about factor income / capital income?
        \item how about all are equally poor?
        \item alternatives: human development index (inequality-adjusted)
        \item income/economic mobility?
        \item how about migration?
    \end{itemize}
    
    }
\end{frame}

\begin{frame}{Elephant Curve}
\mycomment{
    The so-called "elephant curve", also called Lakner-Milanovic-curve, shows the global income distribution. Income has improved for almost everyone; however, upper incomes have profited more. 

    At the right, at very high incomes, there are populations that have not been profiting at all; while a few very rich in income have gained significant wins in the last decades.
}
    \showmybitmap{../shared/bitmaps/elephantcurve.png}    
    \cite{milanovic13}
\end{frame}



\begin{frame}{Is Trade the Culprit?}
    \mycomment{
        Criticism of world trade became prominent in the decades of globalisation after the end of the cold war. 
    }
    \mytext{
    \begin{itemize}
        \item If the farmer can't swim, he blames the swim suit
        \item Lack of standards and global governance
        \item Standards instead of tariffs
        \item Opacity of origin
        \item Supply Chain Law
        \item Brand protection etc.
    \end{itemize}
    }
\end{frame}


% \begin{frame}{Recap Microeconomics}
%     \showmytikz[]{../shared/tikz/tikz-monopoly.tex}
%     \mycomment{
%         In the competition case, more suppliers will enter the market until the \textbf{price equals marginal costs.}  
        
%         Everyone is a \textbf{price taker.} The market determines price and volume. 

%         What's not to love?
%     }
% \end{frame}
% \begin{frame}{Recap Microeconomics}
%     \showmytikz[1]{../shared/tikz/tikz-monopoly.tex}
%     \mycomment{
%         The result is a \textbf{low price,} a \textbf{large quantity,} a \textbf{large consumer surplus} and \textbf{no profits.} Total welfare is as large as it can get.

%         (To be precise: Whether there is a producer surplus depends on the interpretation of the cost curve as composed by individual costs or uniform for all sellers. If it is uniform, there are indeed no profits. If there are some sellers who manage to have lower costs than others, they may make a profit. Then, the sum of all these profits is $PS$.)
%         )
%     }
% \end{frame}
% \begin{frame}{Recap Microeconomics}
%     \showmytikz[1,2]{../shared/tikz/tikz-monopoly.tex}
%     \mycomment{
%         A monopolist can choose any price. And he will set the price so that his profit is maximum. Profit is revenue minus costs:

%         $$ P(x) = R(x) - C(x)  \rightarrow max$$

%         $$ 0 = \frac{dR(x)}{dx} - \frac{dC(x)}{dx} $$

%         $$ 0 = MR - MC$$
%         $$ MR = MC $$
        

%         Profits for the monopolist are at a maximum where \textbf{marginal revenue equals marginal cost.}
        
%         Point $Y$ is called \textbf{Cournot point.}
%         }
% \end{frame}
% \begin{frame}{Recap Microeconomics: Monopoly}
%     \showmytikz[1,2,3]{../shared/tikz/tikz-monopoly.tex}
%     \mycomment{
%         Welfare in the monopoly case is much lower than in the competition case. In particular, consumer surplus is much lower, while a large part of the surplus now is the producer surplus, i.e. profit for the monopolist.

%         What's not to love?

%         Would a tax on the producer surplus and a transfer to the consumer be a good idea? Hint: that would mean extra costs (the tax) for the producer, and more demand (due to increased income) by the consumer...
%     }
% \end{frame}

% \begin{frame}{Recap Microeconomics: Monopoly Regulation}
%     \showmytikz[1,2,3,4]{../shared/tikz/tikz-cournot.tex}
%     \mycomment{
%         It is possible to regulate a monopolist by forcing its price down towards the competitive outcome ("first best"), by imposing a price cap.

%         This is typically above marginal costs to make sure that fixed costs are served. The price cap would be set at price equals average costs ("second best").

%         Note that such a price cap will not work in a competition situation!
%     }
% \end{frame}



\begin{frame}{Intra-Industry Trade}
    \mytext{
    \begin{itemize}
        \item Japan--Germany trade cars in both directions
        \item Contradicts many theoretical predictions
        \item Leontieff paradoxon: trade between similar economies
        \item Linder hypothesis: due to similar demand
        \item Similar to horizontal integration
        \item rather than vertical integration
    \end{itemize}
    }
\end{frame}

\begin{frame}{Upstream and Downstream}
    \mytext{
    \begin{itemize}
        \item upstream: mining $\rightarrow$ processing
        \item downstream: processing $\rightarrow$ consumer
        \item recycling
        \item urban mining
        \item cradle-to cradle, circular economy
    \end{itemize}
    }
\end{frame}
\begin{frame}{Resources and Reserves}
    \showmytikz{../shared/tikz/tikz-reserves-resources}
    \mycomment{
    As resources dry up, prices increase, which makes exploration and extraction of resources economically attractive that were before considered out of reach.

    Thus, there will likely never be a last drop of oil. 

    Some materials - such as fossile fuels - are burnt and then gone. Others, such as metals, can be recycled. Urban mining and recycling will play an increasing role in the future (circular economy).
    }
\end{frame}

